\documentclass[11pt,a4paper]{article}
% Compilation depuis la racine du projet :
% pdflatex -output-directory=output/pdf rapports/synthese/rapport_synthetique.tex
\usepackage[utf8]{inputenc}
\usepackage[T1]{fontenc}
\usepackage[french]{babel}
\usepackage{lmodern}
\usepackage{amsmath,amssymb}
\usepackage[margin=2cm]{geometry}
\usepackage{microtype}
\usepackage{booktabs}
\usepackage{tabularx}
\usepackage{hyperref}
\usepackage{fancyhdr}
\hypersetup{hidelinks,pdftitle={Projet MLSF - Meowdoku : rapport global de synthèse},
  pdfsubject={Synthèse des questions 1 à 6}}
\pagestyle{fancy}
\fancyhf{}
\fancyhead[L]{\small Projet MLSF -- Meowdoku}
\fancyhead[R]{\small Rapport global -- Questions 1 à 6}
\fancyfoot[C]{\thepage\,/\,\pageref{fin-rapport}}
\setlength{\headheight}{14pt}
\setlength{\parindent}{0pt}
\setlength{\parskip}{5pt}

\begin{document}
\begin{center}
  {\LARGE\bfseries Résoudre Meowdoku par SAT}\par\vspace{5pt}
  {\large Rapport global de synthèse}
\end{center}

\textbf{Objectif.} Construire en Python 3.12 une chaîne complète de
résolution de Meowdoku, depuis la modélisation logique jusqu'à un solveur
optimisé, puis évaluer sa correction et ses performances. Le programme
utilise uniquement la bibliothèque standard de Python.

\begin{center}
\fbox{\small Grille $\longrightarrow$ CNF / DIMACS $\longrightarrow$
solveur SAT $\longrightarrow$ placement vérifié}
\end{center}

\section*{1. Modélisation et génération des grilles}

Une variable $x_{i,j}$ représente la présence d'un chat dans la case
$(i,j)$ d'une grille de $N$ lignes, $M$ colonnes et $K$ régions.
L'encodage comporte $NM$ variables et deux formes de clauses :
\[
  \underbrace{\bigvee_{u\in R} x_u}_{\text{au moins un chat dans chaque région }R}
  \qquad\text{et}\qquad
  \underbrace{\neg x_u\lor\neg x_v}_{\text{deux cases incompatibles }u,v}.
\]
Les paires incompatibles imposent au plus un chat par ligne, colonne et
région, ainsi que l'absence de chats voisins, diagonales comprises.
Ainsi, la formule est satisfiable \textbf{si et seulement si} la grille
admet un placement valide. Les clauses identiques sont regroupées et la
numérotation DIMACS est $\operatorname{id}(i,j)=(i-1)M+j$.

Le générateur se limite au cas autorisé $N=M=K=n$. Il construit d'abord
un placement valide par retour arrière, puis fait grandir une région
connexe autour de chaque chat. Une solution existe donc par construction,
mais son unicité n'est pas garantie. Les tailles $n=2$ et $n=3$, impossibles,
sont rejetées.

\section*{2. Lecture DIMACS et DPLL de référence}

Les clauses sont représentées par des ensembles immuables de littéraux
entiers. Les deux versions demandées de DPLL sont implémentées : la
première décide SAT/UNSAT ; la seconde fournit aussi une valuation en cas
de satisfiabilité.

Le solveur détecte une clause vide ou une formule vide, applique la
propagation unitaire et la règle du littéral pur, puis choisit aléatoirement
une variable et sa valeur. En cas d'échec, il explore la valeur opposée.
Cette version sert de référence ; la recopie de la formule à chaque
simplification et la taille de l'arbre de recherche limitent sa vitesse.

\section*{3. Résolution et restitution d'une grille}

\texttt{meowdoku.py} convertit la grille en DIMACS, appelle le DPLL de
référence et transforme les variables vraies en coordonnées de chats.
Il vérifie à la fois les clauses SAT et, indépendamment, les règles du jeu.
Le fichier résultat reprend l'entrée, les coordonnées par région et une
grille marquée de \texttt{*}. Pour une grille UNSAT, aucun résultat n'est
écrit. Le moteur optimisé de la question 6 est accessible séparément par
\texttt{sat.py}.

\newpage
\section*{4. Heuristiques de branchement}

L'algorithme DPLL est conservé ; seule sa fonction de choix devient
paramétrable. Cette séparation permet de comparer les heuristiques avec
les mêmes règles de propagation et de construction du modèle.
Les 16 variantes couvrent les critères suivants :

\begin{center}\small
\begin{tabularx}{\linewidth}{lX}
\toprule
Famille & Critère et justification \\
\midrule
H0 / H1 & Hasard / première variable libre : références peu coûteuses. \\
H2 / H3 & Variable / littéral le plus fréquent : agir sur de nombreuses clauses. \\
H4 / H5 & Fréquences dans les clauses les plus courtes : provoquer des unités
et révéler plus tôt les contradictions. \\
H6 / H7 & Variable / littéral de meilleur poids de Jeroslow--Wang :
$J(l)=\sum_{C\ni l}2^{-|C|}$ privilégie les clauses courtes sans ignorer les autres. \\
H8 & Priorité aux littéraux positifs, départagés par leur poids :
sur Meowdoku, essayer de placer un chat. \\
\bottomrule
\end{tabularx}
\end{center}

H2 à H8 ont une variante \textbf{statique} (scores initiaux) et une
variante \textbf{dynamique} (scores recalculés sur la formule résiduelle).
Le statique économise du calcul, mais son classement peut devenir moins
pertinent. Le dynamique tient compte des simplifications, au prix d'un
parcours supplémentaire de la formule à chaque décision.

H8 répond à la structure de Meowdoku : les nombreuses exclusions binaires
négatives dominent les fréquences, tandis que les clauses positives de
régions sont peu nombreuses. Favoriser un littéral positif déclenche les
exclusions associées à un placement. Ce choix est une hypothèse à tester,
pas une garantie de supériorité sur toute CNF.

\section*{5. Analyse des performances}

\subsection*{5.1 Protocole et lecture des indicateurs}

La campagne porte sur \textbf{173 instances et 2\,768 résolutions} :
50 instances 3-SAT SAT, 50 UNSAT, 23 instances structurées
(17 SAT, 6 UNSAT), 20 Meowdoku SAT de tailles 8 à 16 et 30 grilles
à régions aléatoires (7 SAT, 23 UNSAT). La limite est de 60 s réelles
par résolution ; les processus mesurent chacun leur temps CPU.

Les décisions et retours arrière décrivent l'arbre exploré ; les
propagations et littéraux purs renseignent sur les simplifications.
Le temps CPU mesure leur coût global. Le nombre de résolutions terminées
et le \textbf{PAR-2}, moyenne des temps CPU où un dépassement compte pour
120 s, empêchent d'ignorer les instances difficiles.

Les compteurs et gains ci-dessous sont comparés sur les instances
résolues par \emph{toutes} les heuristiques du groupe, pour ne pas mélanger
travail achevé et compteurs interrompus. Pour les temps, le gain est la
moyenne géométrique des rapports instance par instance,
$G=\exp\!\bigl(\frac1m\sum_i\ln(t_{\mathrm{H0},i}/t_{\mathrm{H},i})\bigr)$.
$G>1$ indique une accélération. Ce sous-ensemble doit être lu avec les
taux de résolution : il exclut certaines instances difficiles.

\begin{center}\small
\begin{tabular}{lrrr}
\toprule
Gain CPU sur H0 & 3-SAT SAT (50) & 3-SAT UNSAT (50) & Meowdoku SAT (19) \\
\midrule
H2s & 6,8 & 7,9 & 0,9 \\
H4s & 7,2 & 7,9 & 0,8 \\
H4d & \textbf{11,0} & \textbf{14,6} & 0,8 \\
H8s & 8,2 & 5,2 & \textbf{3,3} \\
H8d & 7,2 & 6,3 & 2,7 \\
\bottomrule
\end{tabular}
\end{center}
{\small Effectifs communs entre parenthèses. Les cinq variantes résolvent
les 20 grilles Meowdoku SAT.}

\newpage
\subsection*{5.2 Résultats conformes aux intuitions et nuances}

\textbf{Sur le 3-SAT, privilégier les clauses courtes est efficace.}
H4d fait passer la médiane des décisions de 249 à 17 sur les SAT et
d'environ 568 à 40 sur les UNSAT. Ses gains CPU atteignent respectivement
11,0 et 14,6. Les clauses courtes deviennent rapidement unitaires ou
contradictoires : leur priorité réduit l'exploration. Le gain plus marqué
en UNSAT est cohérent avec l'obligation de réfuter toutes les branches
encore possibles ; en SAT, une branche favorable peut suffire.

\textbf{Sur les instances structurées, il faut aussi regarder les délais.}
H4d, H5d et H8d résolvent les 17 SAT, contre 15 pour H0.
Aucune variante ne termine \texttt{public-03} en 60 s. Sur les cinq
UNSAT résolues par toutes, H4s obtient un gain CPU de 25,5 :
ce résultat ne doit pas masquer la sixième instance non résolue.

\textbf{Sur Meowdoku, la polarité est déterminante, mais les décisions
ne suffisent pas à mesurer la difficulté.}
Sur les 19 grilles SAT communes, H8s a une médiane de 10 décisions et
3 retours arrière, contre 54 et 43 pour H0. À l'inverse, H2d atteint
115 décisions mais zéro retour arrière en médiane : exclure les cases
une à une peut être long sans conduire à beaucoup d'impasses.
H8s gagne un facteur 3,3 en temps CPU ; sur les quatre grilles
$16\times16$, sa médiane est de 0,068 s contre 0,530 s pour H0.
Ce rapport de médianes est distinct du gain géométrique global.

Les 23 grilles aléatoires UNSAT sont souvent réfutées après très peu de
décisions. Elles testent la détection de contradictions, mais ne suffisent
pas à représenter les Meowdoku UNSAT difficiles.

\subsection*{5.3 Le recalcul dynamique vaut-il son coût ?}

\begin{center}
\small
\setlength{\tabcolsep}{8pt}
\begin{tabular}{lrrrl}
\toprule
Famille & Décisions s/d & Coût d'un choix d/s & Temps CPU s/d & Résolues s / d \\
\midrule
H2 & 1,28 & 6,3 & 0,87 & 169 / 170 \\
H4 & 1,74 & 5,0 & 1,31 & 171 / 172 \\
H6 & 1,79 & 7,0 & 1,17 & 171 / 171 \\
H8 & 1,34 & 6,4 & 0,97 & 170 / 172 \\
\bottomrule
\end{tabular}
\end{center}
{\small Moyennes géométriques sur les instances résolues par les deux
variantes ; pour les décisions, le rapport est $(D_s+1)/(D_d+1)$.
Les nombres résolus portent sur les 173 instances. Le coût d'un choix
utilise la durée cumulée de l'heuristique divisée par ses décisions.}

Le dynamique fait moins de décisions \emph{en moyenne sur ces paires},
sans garantie instance par instance. Pour H4, le rapport CPU 1,31
correspond à environ 24\,\% de durée en moins ; pour H6, 1,17 correspond
à environ 15\,\%. L'actualisation des petites clauses compense ici son coût.
Pour H2, le rapport 0,87 traduit au contraire environ 15\,\% de durée en
plus : l'amélioration de l'arbre ne rembourse pas les recalculs.

H8 illustre un autre compromis : le dynamique est légèrement plus lent
sur les paires communes, mais résout deux instances supplémentaires.
Sur Meowdoku SAT, H8s consacre environ 5\,\% de sa durée à l'heuristique,
contre 28\,\% pour H8d. Le choix final doit donc associer
\textbf{qualité du branchement, coût du calcul et taux de résolution}.
Ces observations conduisent à conserver une priorité aux placements,
tout en réduisant les parcours de la formule dans le moteur optimisé.

\newpage
\section*{6. Solveur optimisé : conception et résultats}

\subsection*{6.1 Réduire le coût d'une étape de recherche}

Le moteur \texttt{optimized.py} conserve DPLL, mais remplace les copies
de formule par des structures modifiées en place :
\begin{itemize}
  \setlength{\itemsep}{2pt}
  \item Une clause binaire $(a\lor b)$ devient les implications
    $\neg a\Rightarrow b$ et $\neg b\Rightarrow a$. Placer un chat
    propage directement les exclusions des cases incompatibles.
  \item Les clauses longues ont deux littéraux surveillés. Seules les
    clauses dont une surveillance devient fausse doivent être examinées
    pour déplacer la surveillance, propager une unité ou détecter un conflit.
  \item Une pile mémorise les affectations ; le retour arrière annule celles
    du niveau abandonné. La recherche itérative évite une pile d'appels
    récursifs profonde.
\end{itemize}
Les tautologies et doublons sont éliminés au départ. La recherche de
littéraux purs est supprimée pour éviter un parcours global ; cette
règle n'est pas indispensable à la complétude de DPLL.

Pour une CNF composée uniquement de clauses positives et d'exclusions
binaires négatives, le moteur choisit dans la clause positive non satisfaite
qui a le moins de candidats libres. Les scores positifs initiaux de
Jeroslow--Wang départagent ses candidats. Ce choix adapte H8 à Meowdoku :
il actualise la contrainte la plus urgente en ne parcourant que les
clauses positives. Ce n'est pas exactement H8d pour des clauses positives
qui se recouvrent. Les autres CNF utilisent H4d.

\textbf{Correction.} Toute décision garde une branche opposée.
Les implications et surveillances préservent les clauses d'origine.
Au retour arrière, aucune affectation ne devient fausse : elle devient
libre, ce qui permet de conserver les surveillances. Dans la classe de
placement, une fois les clauses positives satisfaites, compléter à faux
les variables restantes préserve les exclusions déjà propagées.
Le moteur reste général ; sa structure n'impose pas qu'une solution existe.

\subsection*{6.2 Une évaluation distincte, proche de l'interface finale}

Le banc de mesure lance \texttt{sat.py} avec son modèle dans un nouveau
processus pour chaque essai. Les 91 instances comprennent les 41 fichiers
publics et les 50 grilles Meowdoku ; les deux modes sont lancés
séquentiellement avec un délai de 60 s, soit 182 essais.
Les statuts attendus proviennent des fichiers publics, de la construction
des grilles SAT et de l'accord observé en Q5 pour les régions aléatoires.

\begin{center}\small
\begin{tabular}{lrr}
\toprule
Indicateur & DPLL initial & Optimisé \\
\midrule
Instances correctement résolues & 89/91 & \textbf{90/91} \\
Benchmarks publics résolus & 28/30 & 29/30 \\
Tests publics d'interface et cas limites & 11/11 & 11/11 \\
Grilles Meowdoku résolues & 50/50 & 50/50 \\
Temps complet sur les 50 grilles Meowdoku & 3,01 s & 1,62 s \\
PAR-2 sur les 91 instances (temps complet) & 3,439 s & 1,421 s \\
\bottomrule
\end{tabular}
\end{center}

Le temps complet inclut Python, les imports, la lecture, la résolution et
la sortie du modèle ; sa vérification intervient ensuite, hors mesure.
Le PAR-2 compte chaque dépassement pour 120 s. Le classement prévu
privilégie d'abord le nombre de réponses correctes, puis leur temps total.
Sur ce corpus, l'optimisation apporte donc un bénéfice de robustesse
en plus du gain de vitesse.

\newpage
\subsection*{6.3 Interpréter les gains et les contre-exemples}

Sur les 50 grilles Meowdoku, le gain moyen géométrique est de
\textbf{4,2 en temps CPU du solveur} et de \textbf{1,5 en temps complet}.
Les coûts fixes du lancement et de la lecture limitent le gain visible
sur les petites instances. Sur les 89 instances résolues par les deux
modes, le gain géométrique en temps complet est de 1,81.

\begin{center}\small
\begin{tabular}{lrrl}
\toprule
Instance publique & DPLL (s) & Optimisé (s) & Observation \\
\midrule
03 & $>60$ & $>60$ & Limite persistante des deux moteurs \\
06 & $>60$ & 0,125 & Une instance supplémentaire résolue \\
12 & 49,840 & 1,306 & Forte réduction du temps de recherche \\
04 & 0,070 & 4,137 & Cas défavorable à l'heuristique finale \\
\bottomrule
\end{tabular}
\end{center}
{\small Temps complets ; valeurs arrondies, avec une seule graine pour DPLL.}

Sur \texttt{public-04}, DPLL trouve un modèle en 13 décisions, contre
50\,088 pour le mode optimisé. Ce contre-exemple montre qu'un branchement
défavorable peut annuler le bénéfice d'une propagation plus efficace.
Il ne suffit pas à conclure que le hasard est préférable en général.

Inversement, sur \texttt{meow\_n16\_s0}, le mode optimisé prend
1\,049 décisions contre 564, mais termine en 0,054 s contre 0,163 s.
Le nombre de décisions seul ne prédit donc pas le temps : leur coût
et celui des propagations ont changé. Le protocole mesure le gain
\emph{combiné} du moteur et de l'heuristique ; il n'isole pas leur
contribution respective.

\subsection*{6.4 Validation et portée des conclusions}

La validation réalisée compte \textbf{54 tests}, dont une comparaison à
un oracle exhaustif sur 2\,000 petites CNF. Les tests couvrent les cas
limites, la propagation, les retours arrière et l'interface. Aucun statut
incorrect ni modèle invalide n'a été observé dans la campagne finale.
Ces contrôles renforcent la confiance, sans remplacer les invariants
de correction.

Les mesures finales utilisent Python 3.12.13 sur macOS 15.7.3 ARM64,
avec une seule exécution et une graine DPLL par instance. Une campagne
à plusieurs graines et répétitions serait nécessaire pour quantifier
la variabilité. Les mesures Q5 et Q6 ne sont pas directement interchangeables :
Q5 compare les heuristiques avec du temps CPU et des processus concurrents ;
Q6 mesure aussi les lancements complets et utilise d'autres graines.

La recherche reste exponentielle, sans apprentissage de clauses, et H4d
parcourt encore la formule à chaque décision. Des tests sur de plus grandes
grilles, des UNSAT plus difficiles et la machine de compétition élargiraient
la validation. Une comparaison à heuristique identique avant/après
changement de propagation isolerait mieux l'origine du gain.

\textbf{Conclusion.} Le projet montre que la performance dépend à la fois
de la structure de l'instance, de la polarité choisie et du coût de la
propagation. L'optimisation répond aux observations de Q5 et améliore
le bilan sur le corpus final, tout en conservant des limites identifiées.

\subsection*{Usage des assistants IA}

Nous avons utilisé des assistants IA lorsque nous rencontrions des blocages
dans la modélisation, afin d'explorer des pistes et de mieux comprendre les
problèmes rencontrés. Ils nous ont également servi de support pour améliorer
la structure générale du code et corriger sa syntaxe. Enfin, nous nous en
sommes servis pour nous aider à rédiger et à reformuler certaines parties de
ce rapport.

\vfill
{\footnotesize\textit{Traçabilité et reproduction.} Sources : rapports détaillés,
\path{results/q5_runs.csv} et \path{results/q6_runs.csv}.
Les mesures existantes sont réanalysées, sans nouvelle campagne.
\par
Q5 : \texttt{uv run python bench\_q5.py report --timeout 60}.
Q6 : \texttt{uv run python bench\_q6.py report}.
\par
Résolution : \texttt{uv run python sat.py instance.cnf --mode optimized --output-model}.}
\label{fin-rapport}
\end{document}
